\chapter{Literature Review and Related Work}
\label{chap:related_work}

This chapter establishes the theoretical and empirical context of this dissertation by reviewing four foundational pillars of literature: machine-generated text detection paradigms, computational morphology for Turkic languages, automated fact-checking and evidence retrieval systems, and contrastive representation learning for out-of-distribution robustness.

\section{Taxonomy of AI-Generated Text Detection}
\label{sec:rel_mgt_detection}

Automated techniques for detecting machine-generated text broadly fall into three principal methodologies: statistical metric-based methods, generation-time watermarking, and supervised neural classifiers.

\subsection{Statistical and Zero-Shot Detectors}
Statistical detectors inspect distributional artifacts within token probability curves generated by language models. The seminal Giant Language Model Test Room (GLTR) \citep{gehrmann2019gltr} visualizes token rank distributions, observing that autoregressive decoders draw heavily from the top-$k$ predictive head compared to human writers. DetectGPT \citep{mitchell2023detectgpt} formalizes this intuition through probability curvature: by applying small masked perturbations to candidate texts, it measures whether the original sample resides in a local maximum of the model's log-probability landscape. Fast-DetectGPT \citep{bao2024fastdetectgpt} extends this concept by analyzing conditional probability curvature directly, achieving substantial speedups. Binoculars \citep{hans2024binoculars} utilizes the ratio between perplexity and cross-perplexity across two disparate foundation models, enabling zero-shot thresholding without training. However, these zero-shot methods assume white-box access to token probability distributions, which is frequently inaccessible for commercial closed APIs or low-resource target languages lacking matching reference LLMs.

\subsection{Watermarking Protocols}
Generation-phase watermarking schemes, initiated by \citet{kirchenbauer2023watermark}, partition vocabulary tokens into dynamic ``green'' and ``red'' lists based on preceding token hashes. During decoding, sampling from the green list is statistically promoted via logit biasing, creating an invisible, verifiable cryptographic trace. While resilient against minor edits, watermarking requires cooperative generator deployments and cannot address unauthorized or open-source foundation models.

\subsection{Supervised Neural Classifiers and Generalization Vulnerabilities}
Supervised neural detectors fine-tune bidirectional transformers (such as BERT \citep{devlin2019bert} and RoBERTa \citep{liu2019roberta}) to classify text representations as human or machine-generated \citep{solaiman2019release}. Although supervised classifiers attain exceptional in-distribution accuracy, recent benchmark studies, including RAID \citep{wang2024raid} and GenAIDetect \citep{genaidetection2025}, demonstrate acute vulnerability to out-of-distribution shifts: detectors trained on one generator family (e.g., GPT-3.5) suffer severe performance degradation when tested against modern alternatives (e.g., LLaMA \citep{touvron2023llama} or Qwen \citep{qwen2024}). Furthermore, \citet{liang2023gptdetectors} revealed that commercial detectors systematically misclassify non-native and low-resource text due to lower lexical perplexity, highlighting an urgent fairness and calibration challenge.

\section{Computational Morphology and NLP for Turkic Languages}
\label{sec:rel_turkic_morphology}

The formal computational modeling of agglutinative morphology traces back to two-level morphology \citep{koskenniemi1983twolevel} and finite-state transducer (FST) theory. \citet{oflazer1994twolevel} developed comprehensive two-level morphological models for Turkish, demonstrating that sequential morphotactics can be modeled through composed finite-state automata. In the Kipchak language family, \citet{tyers2016finite} constructed open-source finite-state transducers for Kazakh, Tatar, and Bashkir within the Apertium platform, establishing foundational phonological rule mappings.

Subsequent investigations have examined morphological segmentation for downstream language processing tasks. \citet{makhambetov2013kazakh,makhambetov2015data} assembled the initial large-scale Kazakh linguistic corpora and developed data-driven statistical stemmers. \citet{bekmanova2021kazakh} combined rule-based affix dictionaries with semantic processing to enhance Kazakh information retrieval. With the advent of subword tokenization algorithms (Byte-Pair Encoding \citep{sennrich2016subword} and SentencePiece \citep{kudo2018sentencepiece}), modern multilingual models (mBERT \citep{devlin2019bert} and XLM-RoBERTa \citep{conneau2020xlmr}) and regional models (KazBERT \citep{eraly2021kazbert} and KazRoBERTa \citep{sagyndyk2025kazroberta}) emerged. However, these tokenizers lack linguistic supervision, fragmenting agglutinated suffixes into irregular token sequences that distort distributional statistics.

\section{Automated Fact Verification and Retrieval Systems}
\label{sec:rel_fact_verification}

Fact verification research was catalyzed by the FEVER (Fact Extraction and VERification) benchmark \citep{thorne2018fever}, which formalized verification as a multi-stage task: selecting relevant evidence sentences from a textual knowledge base (Wikipedia) and determining whether the evidence \textsc{Supports}, \textsc{Refutes}, or provides \textsc{Not Enough Info (NEI)} regarding an asserted claim. Subsequent iterations expanded this paradigm: FEVEROUS \citep{aly2021feverous} introduced verification over structured tables and unstructured text, while SciFact \citep{wadden2020fact} examined scientific claim verification.

\subsection{Lexical Shortcuts and Debiasing}
A central finding in fact verification is the susceptibility of deep classifiers to superficial lexical heuristics \citep{schuster2019fever}. Models frequently exploit claim-only syntactic cues (e.g., negating particles indicating refutation) without genuinely retrieving or conditioning upon evidence. Addressing this challenge requires adversarial benchmark protocols, such as symmetric claim generation and hard neutral perturbations.

\subsection{Dense vs. Hybrid Evidence Retrieval}
Evidence retrieval architectures combine sparse inverted index methods (such as BM25 \citep{robertson2009bm25}) with dense dual-encoders (such as DPR \citep{karpukhin2020dpr} and Contriever \citep{izacard2021contriever}). While dense retrieval maps semantic concepts into shared vector spaces, it frequently stumbles upon rare named entities and low-frequency inflectional variants. Reciprocal Rank Fusion (RRF) \citep{cormack2009reciprocal} offers a robust, parameter-free mechanism for fusing sparse lexical rankings and dense semantic rankings, providing strong empirical foundations for Retrieval-Augmented Generation (RAG) \citep{lewis2020retrieval}. In agglutinative languages, cross-lingual retrieval benchmarks such as TyDi QA \citep{clark2020tydi} demonstrate that dense retrievers alone fail without morphological preprocessing.

\section{Representation Learning, Contrastive Generalization, and Calibration}
\label{sec:rel_representation_learning}

To overcome the brittle domain boundaries of standard cross-entropy optimization, modern deep learning leverages contrastive representation learning. Supervised Contrastive Learning (SupCon) \citep{khosla2020supervised} generalizes self-supervised contrastive objectives by pulling normalized latent representations of samples from the same semantic class together while repelling representations of different classes:
\begin{equation}
    \mathcal{L}_{\text{SupCon}} = \sum_{i \in I} \frac{-1}{|P(i)|} \sum_{p \in P(i)} \log \frac{\exp(\bm{z}_i \cdot \bm{z}_p / \tau)}{\sum_{a \in A(i)} \exp(\bm{z}_i \cdot \bm{z}_a / \tau)}
\end{equation}
SupCon induces hyper-spherical clustering in latent space, which provides superior robustness against domain shifts, noisy labels, and out-of-distribution samples. Furthermore, temperature scaling and Platt scaling \citep{guo2017calibration} ensure that probabilistic outputs reflect true empirical likelihoods, preventing overconfident misclassifications in safety-critical deployments.

\section{Summary of Research Gaps}
\label{sec:rel_gaps}

Existing literature reveals three acute, unaddressed voids:
\begin{enumerate}
    \item \textbf{Absence of Morphologically-Aware AI Detectors:} AI text detection has overlooked agglutinative morphology, causing massive false positive spikes on short colloquial texts due to unmitigated subword token inflation.
    \item \textbf{Lack of Agglutinative Fact-Checking Benchmarks:} Low-resource Turkic languages lack standardized, human-annotated fact verification benchmarks with adversarial controls against lexical shortcuts.
    \item \textbf{Conflation of Machine Origin and Truth:} Existing systems lack a unified mathematical framework that simultaneously quantifies generative likelihood and factual reliability.
\end{enumerate}
The subsequent chapters of this dissertation systematically address each of these identified gaps.
